\chapter{Динамика системы подвижности на кривошипно-шатунных опорах} \label{chapt3}

Трехстепенные электромеханические платформы, рассмотренные в предыдущих главах в силу особенностей выбранного привода и кинематической схемы обладают оптимальными характеристиками с точки зрения энергопотребления, затрат на эксплуатацию и обслуживание.  К основным недостаткам данного типа СП относят сравнительно низкое качество имитации движений, требующих больших ускорений – реакция с запозданием при проезде препятствий, неправдоподобные ощущения при ударах модели техники об объекты на 3D-сцене, пр. Данная особенность как правило связана с невысоким значением перегрузочной способности по вращающему моменту в используемых асинхронных двигателях. Помимо этого, при определенной конфигурации кабины тренажера (высоко расположенный центр тяжести, большая масса кабины), мотор-редукторы СП могут длительное время находиться в пусковом режиме работы. В связи с этим, повышенное внимание при разработке алгоритмов управления СП уделяется точному расчету рабочей зоны как по кинематическим, там и динамическим ограничениям платформы.

Динамический расчёт платформ различной конфигурации достаточно подробно описан в литературе \cite{Sloushch2,2005a,2009a,2010c}, далее приведены выкладки для 3DOF платформы с использованием метода силового расчёта. 

\section{Уравнения обратной задачи динамики}

Схема 3DOF платформы с изображением сил приведена на рис. 1. Особенностью используемой кинематической схемы является центральный карданный шарнир, удерживающий платформу от перемещений по курсу. 
 
\begin{figure}[h]
	\center
	\includegraphics[width=0.75\textwidth]{3dof_kinematics_scheme_dynamic}
	\caption{Кинематическая схема подвижности 3DOF}
	\label{img:3dof_dynamic_scheme}
\end{figure}

Для решения обратной задачи динамики требуется определить силы реакции шатунов на платформу. Т.к. векторы данных сил всегда направлены вдоль шатунов, эти векторы для удобства можно представить в виде произведений $k_1S_1$, $k_2S_2$, $k_3S_3$, где $k_1$, $k_2$, $k_3$ – масштабные коэффициенты, $S_1$, $S_2$, $S_3$ – векторы положения шатунов в неподвижной системе координат, которые определяются по координатам точек крепления шатунов к кривошипам и подвижной платформе.

Обозначим за $S_1^0$, $S_2^0$, $S_3^0$ векторы от центра тяжести до точек крепления шатунов к подвижной платформе. Также ведем обозначения вертикальных проекций векторов $S_{iz}$, нормированных по длине шатуна в земной и подвижной системе координат:
\begin{equation}
\label{eq:DynamicVerticaldPojections}
	F_i = \frac{1}{L}S_{i,z};~ D_i = \frac{1}{L}S_{i,z}^0;~i=1\dots3
\end{equation}

Т.к. движение по курсу отсутствует, при составлении уравнения сил учитывается лишь вертикальная составляющая:
\begin{equation}
\label{eq:DynamicVerticalForcesEq1}
	m\ddot{h} = -k_sh + f_0 - mg + \sum_{i=1}^3 k_iS_{i,z},
\end{equation}
где $k_s$ -- коэффициент упругости пружины, $h$ -- вертикальное смещение, $m$ -- масса подвижной платформы с закрепленным на ней объектом, $f_0$ -- значение подъемной силы, создаваемой пружиной при $h=0$.

Перепишем уравнение относительно нормированных величин $F_i$:
\begin{equation}
\label{eq:DynamicVerticalForcesEq2}
	\sum_{i=1}^3k_iF_i = \frac{1}{L}\left(m\ddot{h}+k_sh-f_0+mg\right)
\end{equation}

Далее запишем уравнение моментов движения твердого тела в связанной системе координат. Уравнение моментов записывается относительно центра тяжести подвижной платформы с закрепленным объектом. Положение центра тяжести в системе координат платформы обозначим $X^0_C = \left[x_c, y_c, z_c\right]^T$. Тогда момент, создаваемый центральной опорой будет равен:
\begin{equation}
\label{eq:TorqueFromCentralSupport}
M_S^0 = -X_C^0 \times F_S^0 = -\left(X_C^0 \times e_S\right)\left(f_0 - k_sh\right),
\end{equation}

где $F^0_S$ – сила, действующая на платформу со стороны центральной опоры, $e_3 = \left[0~0~1\right]^T$ -– единичный вектор, направленный по оси $Z$. В дальнейшем будут использоваться первые две компоненты векторов моментов, т.к. вращение по курсу не рассматривается. Для того, чтобы учесть этот факт, введем вспомогательную матрицу $E_{12}$:
\begin{equation}
\label{eq:AuxiliaryUnitMatrix}
E_{12} = \begin{pmatrix}
1 & 0 & 0 \\
0 & 1 & 0
\end{pmatrix},
\end{equation}

Рассчитаем моменты в плоскости XY от каждого шатуна в системе координат платформы:
\begin{equation}
\label{eq:TorquesFromRodsInMovingCS}
\begin{aligned}
& E_{12}M_1^0 = E_{12}k_1\begin{pmatrix}
-a-x_c \\ b-y_c \\ -z_c
\end{pmatrix} \times S_1^0 = \\
& = k_1\left\{
\left[\begin{pmatrix}b \\ a\end{pmatrix} + 
x_c\begin{pmatrix}0 \\ 1\end{pmatrix} -
y_c\begin{pmatrix}1 \\ 0\end{pmatrix}
\right]S_{1,z}^0+z_c\begin{pmatrix}S_{1,y}^0 \\ -S_{1,x}^0\end{pmatrix}\right\} \\
& E_{12}M_2^0 = E_{12}k_2\begin{pmatrix}
-a-x_c \\ -b-y_c \\ -z_c
\end{pmatrix} \times S_2^0 = \\
& = k_2\left\{
\left[\begin{pmatrix}-b \\ a\end{pmatrix} + 
x_c\begin{pmatrix}0 \\ 1\end{pmatrix} -
y_c\begin{pmatrix}1 \\ 0\end{pmatrix}
\right]S_{2,z}^0+z_c\begin{pmatrix}S_{2,y}^0 \\ -S_{2,x}^0\end{pmatrix}\right\} \\
& E_{12}M_3^0 = E_{12}k_3\begin{pmatrix}
c-x_c \\ -y_c \\ -z_c
\end{pmatrix} \times S_3^0 = \\
& = k_3\left\{
\left[\begin{pmatrix}0 \\ c\end{pmatrix} + 
x_c\begin{pmatrix}0 \\ 1\end{pmatrix} -
y_c\begin{pmatrix}1 \\ 0\end{pmatrix}
\right]S_{3,z}^0+z_c\begin{pmatrix}S_{3,y}^0 \\ -S_{3,x}^0\end{pmatrix}\right\} \\
\end{aligned}
\end{equation}

Запишем угловую скорость вращения платформы в подвижной системе координат. Положение произвольной точки $r$ в неподвижной СК определяется как $r^0 = R_xR_yr = Pr$. Тогда линейная скорость точек платформы $V = -\Omega\times\left(Pr\right) = \dot{P}r$. Равенство верно для любой точки $r$, поэтому справедлива замена $P^{-1}x$. Окончательно получим следующее выражение:

\begin{equation}
\label{eq:AngularSpeedFromLinear}
\Omega \times x = -\dot{P} P^{-1} x ~~\forall x\in\Re^3
\end{equation}

Раскроем правую часть равенства \ref{eq:AngularSpeedFromLinear}:

\begin{equation}
\label{eq:AngularSpeedTurnOperatorDiff}
\begin{aligned}
& \dot{P}P^{-1} = \left(\dot{R_x}R_y\dot{\alpha} + R_x\dot{R_y}\dot{\beta}\right)R_y^{-1}R_x^{-1}=\dot{\alpha}\dot{R_x}R_x^{-1} + \dot{\beta}R_x\dot{R_y}R_y^{-1}R_x^{-1}\\
& \dot{R_x} = \begin{pmatrix}
0 & 0 & 0\\
0 & -\sin{\alpha} & \cos{\alpha}\\
0 & -\cos{\alpha} & -\sin{\alpha}
\end{pmatrix},\\
& \dot{R_x}R_x^{-1} = \begin{pmatrix}
0 & 0 & 0\\
0 & 0 & 1\\
0 & -1 & 0
\end{pmatrix},~
 \dot{R_y}R_y^{-1} = \begin{pmatrix}
0 & 0 & -1\\
0 & 0 & 0\\
1 & 0 & 0
\end{pmatrix}.
\end{aligned}
\end{equation}

Окончательно, равенство \ref{eq:AngularSpeedTurnOperatorDiff2} преобразуется к виду:

\begin{equation}
\label{eq:AngularSpeedTurnOperatorDiff2}
\Omega x = \begin{pmatrix}
0 & -\Omega_z & \Omega_y \\
\Omega_z & 0 & -\Omega_x \\
-\Omega_y & \Omega_x & 0
\end{pmatrix}x = \begin{pmatrix}
0 & -\dot{\beta}\sin{\alpha} & -\dot{\beta}\cos{\alpha} \\
\dot{\beta}\sin{\alpha} & 0 & \dot{\alpha} \\
\dot{\beta}\cos{\alpha} & -\dot{\alpha} & 0
\end{pmatrix}x
\end{equation}

С учетом матрицы $E_{12}$ вектор угловых скоростей имеет вид: 

\begin{equation}
\label{eq:AngularSpeedInMovingCS}
E_{12}\Omega^0 = \begin{pmatrix}
\Omega_x^0 \\ \Omega_y^0 \\ \Omega_z^0
\end{pmatrix} = 
E_{12}\begin{pmatrix}
\dot{\alpha}\cos{\beta} \\ \dot{\beta} \\ -\dot{\alpha}\sin{\beta}
\end{pmatrix} = 
\begin{pmatrix}
\dot{\alpha}\cos{\beta} \\ \dot{\beta}
\end{pmatrix}
\end{equation}

Уравнение моментов с учетом введенной в \ref{eq:AuxiliaryUnitMatrix} матрицы $E_{12}$ имеет вид:
\begin{equation}
\label{eq:TorquesEqCommon}
E_{12}\left[I\dot{\Omega}^0 + \Omega^0 \times \left[I\Omega^0\right] \right] = 
E_{12}\left[M_1^0 + M_2^0 + M_3^0 + M_S^0\right]
\end{equation}

где $I$ – матрица моментов инерции платформы. Т.к. момент инерции записывается относительно центра тяжести, преобладающими являются диагональные компоненты:
\begin{equation}
I = \begin{pmatrix}
I_x & 0 & 0 \\ 0 & I_y & 0 \\ 0 & 0 & I_z
\end{pmatrix}
\end{equation}
Раскроем левую инерционную часть уравнения моментов \ref{eq:TorquesEqCommon}:
\begin{equation}
\label{eq:TorquesEqLeftInertialPart}
E_{12}\left[I\dot{\Omega}^0 + \Omega^0 \times \left[I\Omega^0\right] \right] = \begin{pmatrix}I_x\ddot{\alpha}\cos{\beta} \\ I_y\ddot{\beta}\end{pmatrix} + 
\begin{pmatrix}\left(I_x-I_y-I_z\right)\dot{\beta} \\ \left(I_z-I_x\right)\dot{\alpha}\cos{\beta}\end{pmatrix}\dot{\alpha}\sin{\beta}
\end{equation}

Обозначим инерционные моменты $m_{\alpha}$ и $m_\beta$ для сокращения дальнейших выкладок:
\begin{equation}
\label{eq:TorquesInertial}
\begin{aligned}
& m_{\alpha} = I_x\ddot{\alpha}\cos{\beta} + \left(I_x-I_y-I_z\right)\dot{\alpha}\dot{\beta}\sin{\beta} \\
& m_{\beta} = I_y\ddot{\beta} + \left(I_z-I_x\right)\dot{\alpha}^2\sin{\beta}\cos{\beta}
\end{aligned}
\end{equation}

После подстановки \ref{eq:TorqueFromCentralSupport}, \ref{eq:TorquesFromRodsInMovingCS} и \ref{eq:TorquesInertial} в \ref{eq:TorquesEqCommon}, совместно с \ref{eq:DynamicVerticalForcesEq2} получим систему алгебраических уравнений относительно неизвестных коэффициентов $k_1$, $k_2$, $k_3$, которые в данном случае имеют смысл модулей векторов сил, действующих на шатуны:
\begin{equation}
\label{eq:KoefsLinearSystem1}
\left\{
\begin{aligned}
& F_1k_1 + F_2k_2 + F_3k_3 = \frac{1}{L}\left(m\ddot{h}+k_sh-f_0+mg\right), \\
& b\left(D_1k_1-D_2k_2\right)-Dky_c + S_y^0kz_c = \frac{1}{L}\left(m_{\alpha} + y_c\left(f_0 - k_sh\right)\right), \\
& a\left(D_1k_1+D_2k_2\right)-cD_3k_3+Dkx_c-S_y^0kz_c = \frac{1}{L}\left(m_{\beta} - x_c\left(f_0 - k_sh\right)\right),
\end{aligned}
\right.
\end{equation}

где $D=\left[D_1~D_2~D_3\right],~F=\left[F_1~F_2~F_3\right], k=\left[k_1~k_2~k_3\right]^T,~S^0_x=\left(S^0_{1,x},~S^0_{2,x},~S^0_{3,x}\right)/L, S^0_y=\left(S^0_{1,y},~S^0_{2,y},~S^0_{3,y}\right)/L$.

Прибавим ко второму уравнению в \ref{eq:KoefsLinearSystem1} правую часть первого, умноженную на $y_c$, к третьему –-- ту же правую часть, умноженную на $–x$. При этом для левой части второго и третьего уравнений заменим данное слагаемое на $Fk$ согласно равенству в первом уравнении. Таким образом преобразованные второе и третье уравнения системы \ref{eq:KoefsLinearSystem1}:
\begin{equation}
\label{eq:KoefsLinearSystemSimplification}
\begin{aligned}
& b\left(D_1k_1-D_2k_2\right)-Dky_c + S_y^0kz_c + Fky_c = \frac{1}{L}\left(m_{\alpha} + y_c\left(f_0 - k_sh\right)\right) + \\ & \frac{1}{L}\left(m\ddot{h}+k_sh-f_0+mg\right)y_c, \\
& a\left(D_1k_1+D_2k_2\right)-cD_3k_3+Dkx_c-S_y^0kz_c - Fkx_c = \\ & = \frac{1}{L}\left(m_{\beta} - x_c\left(f_0 - k_sh\right)\right) - \frac{1}{L}\left(m\ddot{h}+k_sh-f_0+mg\right)x_c,
\end{aligned}
\end{equation}

В \ref{eq:KoefsLinearSystemSimplification} выделим $E=D-F$ –- вспомогательную переменную, характеризующую горизонтальные смещения центра тяжести.

Упростив уравнение \ref{eq:KoefsLinearSystemSimplification}, вместе с первым уравнением  \ref{eq:KoefsLinearSystem1} получим преобразованную систему линейных уравнений относительно $k$:

\begin{equation}
\label{eq:KoefsLinearSystem2}
\begin{aligned}
& \begin{pmatrix}
F_1 & F_2 & F_3 \\ 
bD_1 & -bD_2 & 0 \\
aD_1 & aD_2 & -cD_3
\end{pmatrix}k + 
\begin{pmatrix}0 \\ S_y^0 \\ -S_x^0\end{pmatrix}kz_c +
\begin{pmatrix}0 \\ -y_c \\ x_c \end{pmatrix}Ek = \frac{1}{L} 
\begin{pmatrix}
m\ddot{h} \\
m_{\alpha} \\
m_{\beta}	
\end{pmatrix} + \\ 
& + \frac{1}{L}\begin{pmatrix}
k_sh - f_0 + mg \\
y_c(m\ddot{h} + mg) \\
-x_c(m\ddot{h} + mg)
\end{pmatrix}
\end{aligned}
\end{equation}

Далее расчет коэффициентов $k$ осуществляется численно методом простых итераций.

Окончательно, получим формулы для определения вращающих моментов, составив уравнения для каждого кривошипа:
\begin{equation}
\label{eq:TorquesFromCrankAngles}
Q_i  = \begin{pmatrix}
0 \\ R\sin{q_i} \\ -R\cos{q_i}
\end{pmatrix} \times
\left(k_iS_i\right)
\end{equation}

После подстановки $S_i$ в \ref{eq:TorquesFromCrankAngles} получим формулы для моментов относительно рассчитанных значений $k_i$ :
\begin{equation}
\label{eq:TorquesFromCrankAnglesDetailed}
\begin{aligned}
& Q_i = k_iG_i,~i=1,\dots,3 \\
& G_1 = R \left\{ L\sin{q_1}-R\sin{q_1}\cos{q_1}-R^2 \left[ \frac{1}{8b} {\left( \cos{q_1}-\cos{q_2} \right)}^2 \cos{q_1} + \frac{1}{2L}\sin^3{q_1} \right] \right. +\\ 
& \left.+ \frac{R^3}{8Lb} \left( \cos{q_1}-\cos{q_2} \right)^2 \left(\cos{2q_2}+\cos{q_1}\cos{q_2}\right) 
\right\} \\
& G_2 = R \left\{ L\sin{q_2}-R\sin{q_2}\cos{q_2}-R^2 \left[ \frac{2}{8b} {\left( \cos{q_2}-\cos{q_2} \right)}^2 \cos{q_2} + \frac{2}{2L}\sin^3{q_2} \right] \right. +\\ 
& \left.+ \frac{R^3}{8Lb} \left( \cos{q_2}-\cos{q_2} \right)^2 \left(\cos{2q_2}+\cos{q_2}\cos{q_2}\right) 
\right\} \\
& G_3 = R\left\{L\sin{q_3}-R\sin{q_3}\cos{q_3}-R^2\frac{1}{2L}\sin^3q_3\right\}
\end{aligned}
\end{equation}

\section{Исследование решения линейной системы относительно неизвестных сил реакции шатунов}

\subsection{Обращение матрицы линейной системы}
Матрица системы \ref{eq:LinearSystemStatic}:
\begin{equation}
\begin{aligned}
	A_0 = 
	\begin{pmatrix}
	F_1 & F_2 & F_3 \\
	bD_1 & -bD_2 & 0 \\
	aD_1 & aD_2 & -cD_3
	\end{pmatrix}
\end{aligned}
\end{equation}

обращается явно:
\begin{equation}
\begin{aligned}
	A_0^{-1} = \frac{1}{\Delta_0}\begin{pmatrix}
	bcD_2D_3 & cD_3F_2 + aD_2F_3 & bD_2F_3 \\
	bcD_1D_3 & -aD_1F_3-cD_3F_1 & bD_1F_3 \\
	2abD_1D_2 & a\left(D_1F_2-D_2F_1\right) & -b\left(D_1F_2 + D_2F_1\right)
\end{pmatrix}	 
\end{aligned}
\end{equation}

где определитель матрицы $A_0$ равен
\begin{equation}
	\Delta_0 = 2abD_1D_2F_3 + bcD_3\left(D_1F_2 + D_2F_1\right)
\end{equation}

Для решения системы применим метод простой итерации. Запишем уравнения в виде:

\begin{equation}
\begin{aligned}
	A_0k + Hk = B, ~ H = \begin{pmatrix}0 \\ S_y^0 \\ -S_x^0 \end{pmatrix}z_c + \begin{pmatrix}0 \\ -y_c \\ x_c \end{pmatrix}E
\end{aligned}
\end{equation}

Пусть спектральный радиус $\rho\left(A_0^{-1}H\right) < 1$. Тогда решение можно представить в виде ряда:

\begin{equation}
\begin{aligned}
	k & = \left(A_0 + H\right)^{-1}B = A_0^{-1}\left(I + HA_0^{-1}\right)^{-1}B \\ 
	&= A_0^{-1}\sum_{i=0}^{\inf}\left(-HA_0^{-1}\right)^i B = A_0^{-1}\left(\sum_{i=0}^{\inf}C^i\right)B,
\end{aligned}
\end{equation}

где 

\begin{equation}
\begin{aligned}
	C = -H A_0^{-1} = \begin{pmatrix}0 \\ S_y^0 \\ -S_x^0 \end{pmatrix}A_0^{-1}z_c + \begin{pmatrix}0 \\ -y_c \\ x_c \end{pmatrix}EA_0^{-1}
\end{aligned}
\end{equation}

Второе слагаемое мало, т.к. величина $E$ пропорциональна $\kappa^2$, а смещение $x_c,~y_c$ составляют сотые доли метра. Первое слагаемое рассчитывается в следующем разделе.

Особый интерес представляет аппроксимация ряда частичной суммой первого или второго порядка:

\begin{equation}
\begin{aligned}
	\hat{k}^{(1)} = A_0^{-1}B + A_0^{-1}CB, ~~ \hat{k}^{(2)} = A_0^{-1}B + A_0^{-1}CB  + A_0^{-1}C^2B
\end{aligned}
\end{equation}

Запишем матрицу $C$ в виде:

\begin{equation}
\begin{aligned}
	C = -C_0z_c - \begin{pmatrix}0 \\ -y_c \\ x_c \end{pmatrix}EA_0^{-1}, ~C_0 = \begin{pmatrix}0 \\ S_y^0 \\ -S_x^0 \end{pmatrix}A_0^{-1}
\end{aligned}
\end{equation}

Назовем $C$ матрицей итераций.

\subsection{Расчет матрицы итераций}

Введем обозначение для нормированного определителя матрицы $A_0$:
\begin{equation}
\begin{aligned}
	\Delta = \Delta/b = 2aD_1D_2F_3 + cD_3\left(D_1F_2 + D_2F_1\right)
\end{aligned}
\end{equation}

и для столбцов нормированной присоединенной матрицы $A_1 = \Delta A_0^{-1}e_1,~A_2 = \Delta A_0^{-1}e_2,~A_3 = \Delta A_0^{-1}e_3$:

\begin{equation}
\begin{aligned}
	A_1 = \begin{pmatrix}cD_2D_3 \\ cD_1D_3 \\ 2aD_1D_2\end{pmatrix},~ A_2 = \frac{1}{b}\begin{pmatrix}cD_3F_2 + aD_2F_3 \\ -aD_1F_3-cD_3F_1 \\ a\left(D_1F_2-D_2F_1\right)\end{pmatrix}, ~A_3=\begin{pmatrix}D_2F_3 \\ D_1F_3 \\ -\left(D_1F_2 + D_2F_1\right) \end{pmatrix}
\end{aligned}
\end{equation}

Вычислим каждый столбец в матрице:

\begin{equation}
\begin{aligned}
	C_0 = \begin{pmatrix}0 \\ S_y^0 -S_x^0\end{pmatrix}A_0^{-1} = \frac{1}{\Delta}\begin{pmatrix}
	0 & 0 & 0 \\ S_y^0A_1 & S_y^0A_2 & S_y^0A_3 \\ -S_x^0A_1 & -S_x^0A_2 & -S_x^0A_3
	\end{pmatrix}
\end{aligned}
\end{equation}

С точностью до третьего порядка:

\begin{equation*}
	\begin{aligned}
		\frac{1}{\Delta}\begin{pmatrix}S_y^0A_1 \\ -S_x^0\end{pmatrix} & = \frac{2\left(a+c\right)}{\Delta}\left\lbrace\begin{pmatrix}\sin\alpha\cos\beta \\ \sin\beta\end{pmatrix}- \right.\\
		& \left. -\kappa\begin{pmatrix}1\\0\end{pmatrix}\frac{2a\sin{q_3}+c\sin{q_1}+c\sin{q_2}}{2\left(a+c\right)}\left(\cos\alpha-\kappa^2C_{1,\beta}^2\right) - \right.\\
		& \left. -\frac{1}{2}\kappa^2 \begin{pmatrix} \sin{\alpha} \\ \sin{\beta} \end{pmatrix} \frac{2a\sin^2{q_3}+c\left(\sin^2{q_1}+\sin^2{q_2}\right)}{2\left(a+c\right)}-\kappa^2 \begin{pmatrix} \sin{\alpha} \\ 0 \end{pmatrix} C^2_{1,\beta} + \right.\\
		& \left. \frac{1}{2} \kappa^3 \begin{pmatrix} 1 \\ 0 \end{pmatrix} \frac{2asq_3 \left(sq_1^2+sq_2^2 \right) + csq_1\left( sq_2^2+sq_3^2 \right) + csq_2\left(sq_1^2+sq_3^2\right)}{2\left(a+c\right)} \right\rbrace\\
		 S_y^0 A_2 & \approx -\kappa \frac{a+c}{b} \left( \sin{q_1} - \sin{q_2} \right) \left\lbrace 1 + \frac{1}{2} \kappa^2 \left(4\lambda - C^2_{1,\beta} + \sin{q_1} \sin{q_2} - \sin^2{q_3}\right) \right\rbrace -\\
		& -\kappa^2\frac{a+c}{L}C_{1,\alpha}\\
		S_y^0 A_3 & \approx \kappa\left(2\sin{q_3} - \sin{q_1} - \sin{q_2}\right) \left(1 - \frac{1}{2}\kappa^2\left(C^2_{1,\alpha} + C^2_{1,\beta}\right)\right) + \\
		& + \frac{1}{2}\kappa^3 \left\lbrace \sin{q_1}\left(\sin^2{q_3} + sq^2_2\right) + \sin{q_2}\left(\sin^2{q_3} + sq^2_1\right)\right.\\
		&\left. -\sin{q_3}\left(\sin^2{q_1} + \sin^2{q_2} + sq^2_1 + sq^2_2\right)\right\rbrace
		S^0_x A_2 & \approx \frac{1}{2b}\kappa^3 C_{1,\beta}c\left(\sin^2{q_2} - \sin^2{q_1}\right) \\
		S^0_x A_3 & \approx \frac{1}{2}\kappa^3 C_{1,\beta}\left(2\sin^2{q_3} - \sin^2{q_1} - \sin^2{q_2}\right) + \kappa^2\frac{a+c}{L}C^2_{1,\beta} + \\
		& + \kappa^2\frac{a+c}{L}\left(C^2_{1,\beta} + 2\kappa C_{1,\beta}C_{2,\beta}\right)
	\end{aligned} 
\end{equation*}

Т.к. $\kappa^3$ является малой величиной и умножается на линейную комбинацию тригонометрических функций, компонентами матрицы итераций, содержащими $\kappa^3$ можно пренебречь. Запишем матрицу итераций до второго порядка $\kappa$:

\begin{equation}
\begin{aligned}
	C_0 & \approx \frac{2\left(a+c\right)}{\Delta}\left\lbrace \begin{pmatrix}
	0 & 0 & 0 \\ \sin{\alpha}\cos{\beta} & 0 & 0 \\ \sin{\beta} & 0 & 0
	\end{pmatrix} + \right.\\
	&\left. + k\begin{pmatrix}
	0 & 0 & 0 \\
	-\frac{2a\sin{q_3}+c\sin{q_1}+c\sin{q_2}}{2\left(a+c\right)} & \frac{\sin{q_2}-\sin{q_1}}{2b} & \frac{2\sin{q_3} - \sin{q_1} - \sin{q_2}}{2\left(a+c\right)}\\	
	0 & 0 & 0 \\
	\end{pmatrix} - \right.\\
	& \left.\frac{1}{2L}\kappa^2\begin{pmatrix}
	0 & 0 & 0 \\
	0 & C_{1,\alpha} & 0 \\
	0 & 0 & C^2_{1,\beta}
	\end{pmatrix}\right\rbrace
\end{aligned}
\end{equation}
 Определитель $\Delta$ отличается от $2\left(a+c\right)$ на величину порядка $\kappa^2$. При возведении в квадрат матрицы $z_cC_0$ умножаются на 0 все элементы первого столбца. Множитель $\kappa^2$ останется только в центральном элементе:
 
\begin{equation}
\begin{aligned}
	\kappa^2\left(\frac{2\left(a+c\right)}{\Delta}\right)^2 \frac{z^2_c}{\left(2b\right)^2}\left(\sin{q_2} - \sin{q_1}\right)
\end{aligned}
\end{equation}

которым, как правило можно пренебречь.
	
Поэтому одна итерация в решении системы линейных уравнений дает неплохую аппроксимацию решения.

\subsection{Слагаемое с горизонтальным смещением центра тяжести}

Второе слагаемое в матрице итераций $C$ содержит произведение вектора из смещений $x_cy_c$ на коэффициент $EA_0^{-1}$, где $E = D - F$.

Оценим компоненты строки $E = \left(E_1,~E_2,~E_3\right)$. При $i=1,2,3$:

\begin{equation}
\begin{aligned}
	E_i & \approx \frac{1}{2}\kappa^2 \left(-q^2_i - C^2_{1,\beta} + \sin^2_{q_i}\right) + \kappa^3 \left\lbrace C_{1,\beta}\gamma_i + sq_i\delta_i + \frac{b_i}{2L}\sin{q_i}C^2_{1,\alpha} \right\rbrace \\
	& + \kappa^2\left[C_{1,\alpha}\sin{q_i} - \frac{1}{2}\left(C^2_{1,\alpha} + C^2_{1,\beta}\right)\right] + \\
	& + \kappa^3\left\lbrace C_{2,\alpha}\sin{q_i} + \frac{a}{2L}C^3_{1,\beta} - \left(C_{1,\beta}C_{2,\beta} + C_{1,\beta}C_{2,\alpha}\right)\right\rbrace
\end{aligned}
\end{equation},

где $b_i$ обозначает следующее: $b_1 = b,~b_2 = -b,~0$. Компоненты строки $E$ имеют второй порядок относительно величины $\kappa$.

Оценим каждую компоненту коэффициента $EA^{-1}_0$. Первая компонента:
\begin{equation}
\begin{aligned}
EA^{-1}_0e_1 \approx & \kappa^2\left[\left(C_{1,\alpha} + \kappa C_{2,\alpha} - \kappa^2 C^2_{1,\beta}\right)\frac{2a\sin{q_3} + c\left(\sin{q_1} + \sin{q_2}\right)}{\Delta} - \right.\\
&\left. - \frac{a+c}{\Delta}\left(C^2_{1,\alpha} + C^2_{1,\beta}\right)\left(1 - \kappa^2 C^2_{1,\beta}\right)\right] + \\
& + \kappa^3\frac{2\left(a+c\right)}{\Delta} \left[\frac{a}{2L}C^3_{1,\beta} - \left(C_{1,\beta}C_{2,\beta} + C_{1,\alpha}C_{2,\alpha}\right)\right] + \\
& + \frac{1}{2}\kappa^4\left[ \frac{1}{2}\left(C^2_{1,\alpha} + C^2_{1,\beta}\right) \frac{2a\left(sq^2_1 + sq^2_2\right) + c\left(q^2_1 + q^2_2 + 2q^2_3\right)}{\Delta}\right. \\
& -\left. \frac{2a\sin{q_3}\left(q^2_1 + q^2_2\right) + c\sin{q_1}\left(q^2_2 + q^2_3\right) + c\sin{q_2}\left(q^2_1 + q^2_3\right)}{\Delta} \right]
\end{aligned}
\end{equation}

Вторая компонента:
\begin{equation}
\begin{aligned}
EA^{-1}_0 e_2 & = \frac{1}{\Delta}A_2 = \frac{a+c}{\Delta b}D_3\left(D_1F_2 - D_2F_1\right) \approx\\
& \approx \frac{2\left(a+c\right)}{\Delta}\kappa^2D_3C_{1,\alpha}\left\lbrace \frac{\sin{q_1} - \sin{q_2}}{2b}\left(1 - \frac{1}{2}\kappa C_{2,\alpha}\right) + \kappa \frac{1}{4L}\left(2C^2_{1,\beta} + C^2_{1,\alpha}\right)\right\rbrace.
\end{aligned}
\end{equation}

Третья компонента:
\begin{equation}
\begin{aligned}
EA^{-1}_0 e_3 & \approx -\kappa\sin{\alpha}\cos^2{\beta}\frac{2\sin{q_3} - \sin{q_1} - \sin{q_2}}{\Delta} - \\
& -\frac{1}{2\Delta}\kappa^4 \left\lbrace -\frac{1}{2}C^2_{1,\beta}\left(2\sin^2{q_3} - \sin^2_{q_1} - \sin^2{q_2}\right) \right. \\
& \left. -C^2_{1,\alpha}\left(\sin{q_3} - \sin{q_1}\right)\left(\sin{q_3} - \sin{q_2}\right)\right. \\
& \left. +C_{1,\alpha}\left(\sin{q_1} - \sin{q_3}\right)\left(sq^2_2 - \sin{q_1}\sin{q_3}\right)\right. \\
& \left. -C_{1,\alpha}\left(\sin{q_2} - \sin{q_3}\right)\left(sq^2_1 - \sin{q_2}\sin{q_3}\right)\right\rbrace
\end{aligned}
\end{equation}


\section{Оценка параметров уравнений обратной задачи динамики}
Практически все параметры в (10) определить на практике с достаточной точностью не представляется возможным из-за сложностей в организации процесса измерения. В связи с этим, целесообразно произвести параметрическую идентификацию, используя данные с датчиков, установленных на СП.

\subsection{Идентификация параметров в статике}
Параметрами, в большей степени влияющими на результат расчета движущих моментов являются вес закрепленного на платформе объекта, коэффициент жесткости $k_s$, начальное натяжение пружины $f_0$ и смещение центра тяжести $(x_c, y_c, z_c)$. Эти величины не меняются в течение длительного времени, поэтому целесообразно произвести идентификацию объекта на основе данных эксперимента.
	Данные параметры можно определить в статическом режиме, т.е. измерив вращающие моменты в нескольких стационарных положениях платформы. При этом уравнения системы (11) путем обнуления производных принимают вид:
\begin{equation}
\label{eq:LinearSystemStatic}
\begin{aligned}
& \begin{pmatrix}
\bar{F_1} & \bar{F_2} & \bar{F_3} \\ 
\bar{b}\bar{D_1} & -\bar{b}\bar{D_2} & 0 \\
\bar{a}\bar{D_1} & \bar{a}\bar{D_2} & -\bar{c}\bar{D_3}
\end{pmatrix}\bar{k} + 
\begin{pmatrix}0 \\ \bar{S}_y^0z_c -y_c\bar{E} \\ -\bar{S}_x^0z_c+x_c\bar{E}\end{pmatrix}\bar{k} = \frac{1}{L} 
\begin{pmatrix}
k_s\bar{h} - f_0 + mg \\
y_cmg \\-x_cmg
\end{pmatrix}
\end{aligned}
\end{equation}

Переменные с чертой в \ref{eq:LinearSystemStatic} обозначают константы, рассчитанные по данным эксперимента или полученные из исходных данных, без черты – идентифицируемые параметры \cite{Graupe}.

Введем дополнительные расчетные величины, которые в \ref{eq:LinearSystemStatic} не зависят от идентифицируемых параметров:
\begin{equation}
\label{eq:AdditionalEstimatedVariables}
\begin{pmatrix}
\bar{p_1} \\ \bar{p_2} \\ \bar{p_3} 
\end{pmatrix} = 
\begin{pmatrix}
\bar{F_1} & \bar{F_2} & \bar{F_3} \\ 
\bar{b}\bar{D_1} & -\bar{b}\bar{D_2} & 0 \\
\bar{a}\bar{D_1} & \bar{a}\bar{D_2} & -\bar{c}\bar{D_3}
\end{pmatrix}\bar{k} 
\end{equation}

Тогда после подстановки \ref{eq:AdditionalEstimatedVariables} в \ref{eq:LinearSystemStatic} получим:
\begin{equation}
\label{eq:AdditionalParamsRightPart}
\begin{aligned}
& \bar{p_1} = \frac{1}{L}\left(mg - f_0 + k_sh\right)\\
& \bar{p_2} = \left(\frac{mg}{L}+\bar{E}\bar{k}\right)y_c - \bar{S}^0_y\bar{k}z_c \\
& \bar{p_3} = \left(\frac{mg}{L}+\bar{E}\bar{k}\right)x_c + \bar{S}^0_x\bar{k}z_c
\end{aligned}
\end{equation}

Пусть по результатам эксперимента были получены n стационарных положений платформы. Т.к. в процессе эксперимента вращающие моменты $Q_1,~Q_2,~Q_3$ измеряются, а $G_1,~G_2,~G_3$ рассчитываются по кинематическим уравнениям, вычисляется вектор $k$ путем решения системы (11) методом простых итераций. Далее на каждом из участков по формуле (18) вычисляются значения  $p_{1,i},~p_{2,i},~p_{3,i},~i = 1,\dots,n$.

Вместо величин $\left(x_c, y_c\right)$ определяются сразу произведения $\left(x_c, y_c\right) = \left(x_c, y_c\right) mg/L$.

Для определения параметров используется метод наименьших квадратов с дополнительными весовыми множителями [3]:

\begin{equation}
\label{eq:LeastSquareMethodFormulation}
\begin{aligned}
& J_c\left(x_c,y_c,z_c\right)=\sum_{i=1}^n\left[
c_{i,x}\left|\bar{p}_{3,i}-\left(\bar{S}_x^0\bar{k}\right)_iz_c + \bar{x_c}\right|^2 + c_{i,y}\left|\bar{p}_{2,i}-\left(\bar{S}_y^0\bar{k}\right)_iz_c  \bar{y_c}\right|^2\right], \\
& J_c\left(k_s,~f_0\right)=\sum_{i=1}^n{\left|\bar{p}_{1,i}-k_s\bar{h_i}-\frac{mg-f_0}{L}\right|}^2, \\
& J_c\left(x_c,~y_c,~z_c\right)\rightarrow\min, \\ 
& J_c\left(k_s,~f_0\right)\rightarrow\min.
\end{aligned}
\end{equation}

где $c_{i,x}$ и $c_{i,y}$  – весовые множители, $i = 1,\dots,n$ – номера стационарных участков.

\subsection{Идентификация в режиме гармонических колебаний}
Предположим, что в записи имеются несколько участков, на которых углы крена и тангажа меняются периодически. Экспериментально установлено, что среди инерциальных моментов сил только чистые вторые производные от углов крена и тангажа могут принимать значения, влияющие на общий суммарный момент сил. Это значит, что:

\begin{equation}
\begin{aligned}
m_{s,\alpha} & \approx 0,\\
m_{s,\beta} & \approx 0,\\
\left(I_y - I_x - I_z\right)\dot{\alpha}\dot{\beta}\sin{\beta} & \approx 0,\\
\left(I_z - I_x\right)\dot{\alpha}^2 \sin{\beta}\cos{\beta} & \approx 0.
\end{aligned}
\end{equation}

Наиболее значимый вклад вносит компонента $\begin{pmatrix}0 & I_x\ddot{\alpha}\cos{\beta} & I_y \ddot{\beta} \end{pmatrix}^T$. 

Поскольку смещения центра тяжести и параметры сил оценены по стационарным участкам, в каждый момент времени коэффициенты сил $k$ можно представить в виде:
\begin{equation}
\begin{aligned}
	k = k_0 + k_{\alpha} + k_{\beta}
\end{aligned}
\end{equation}

где 
$k_0 = A^{-1}B,~k_{\alpha}=A^{-1}\frac{1}{L}\begin{pmatrix}0 \\ I_x\ddot{\alpha}\cos{\beta}\end{pmatrix},~k_{\beta}=A^{-1}\frac{1}{L}\begin{pmatrix}0 \\ 0 \\ I_y\ddot{\alpha}\cos{\beta}\end{pmatrix},~$

и

\begin{equation}
\begin{aligned}
	A & = \begin{pmatrix}F_1 & F_2 & F_3 \\ bD_1 & -bD_2 & 0 \\ aD_1 & aD_2 & -cD_2\end{pmatrix}+
	\begin{pmatrix}0 \\ S_y^0 \\ -S_x^0\end{pmatrix}z_c + \begin{pmatrix}0 \\ -y_c \\ x_c\end{pmatrix}E, \\
	B &= \begin{pmatrix}
	\frac{1}{L}m\ddot{h} + \frac{1}{L}k_sh + \bar{f}\\
	\bar{y_c}\left(\ddot{h}/g + 1\right) \\
	-\bar{x_c}\left(\ddot{h}/g + 1\right)
	\end{pmatrix}
\end{aligned}
\end{equation}

Оценка моментов сил на валах кривошипов раскладывается на три компоненты:

\begin{equation}
\begin{aligned}
	m = m_0 + m_{\alpha} + m_{\beta},~m_0 = Gk_0,~m_{\alpha} = Gk_{\alpha},~m_{\beta} = Gk_{\beta}
\end{aligned}
\end{equation}

На каждом однородном участке колебаний с постоянной частотой это простые гармоники. При помощи БПФ определяется частота $F$ простых гармоник вращающих моментов. Величина $F$ измеряется в количестве периодов на длину выбранного участка сигнала $N$.
Далее определяются комплексные амплитуды этих гармоник:

\begin{equation}
\begin{aligned}
	m_0(t) = \operatorname{Re}\left(A_0 \exp{\frac{2\pi i}{N}Ft}\right),\\
	m_{\alpha}(t) = \operatorname{Re}\left(A_{\alpha} \exp{\frac{2\pi i}{N}Ft}\right),\\
	m_{\beta}(t) = \operatorname{Re}\left(A_{\beta} \exp{\frac{2\pi i}{N}Ft}\right).
\end{aligned}
\end{equation}

Каждая амплитуда $A_0, A_{\alpha}, A_{\beta}$ имеет 3 компоненты, соответствующие кривошипам. Аналогично выделяется первая гармоника измеряемой функции моментов на валах $m_{obs}$. Вектор её амплитуд обозначим $A_{obs}$. Неточно известны множители, с которыми линейные и угловые ускорения влияют на моменты на валах. Это масса кабины m и её моменты инерции $I_x$ и $I_y$. Однако влияние линейных ускорений $h$ на порядок меньше, чем угловых ускорений. Поэтому вместо $m$ можно подставить приблизительную оценку этой величины. Рассмотрим задачу настройки параметров $I_x$ и $I_y$. В измерениях $m_{obs}(t)$ возможны запаздывания, которые проявляются в изменении фазы в комплексной амплитуде $A_{obs}$. Поскольку смещение фазы неизвестно, то настройку будем проводить только по вещественным амплитудам гармоник. Пусть $i = 1,\ldots,n$ — номера участков сигнала с гармоническим изменением углов кривошипов. Определим функционал качества:

\begin{equation}
\begin{aligned}
	J\left(c_x,c_y\right) = \sum^n_{i=1}\left(\|A_{obs,i}\|^2 - \|A_{0,i} + c_xA_{\alpha,i} + c_yA_{\beta,i}\|^2\right)^2
\end{aligned}
\end{equation}

Этот функционал минимизируется по $\left(c_x, c_y\right)$ градиентным методом. При расчете предварительно задаются границы $c_{\min},~c_{\max}$ допустимых отклонений коэффициентов от 1.

После того, как коэффициенты настроены, получаем оценки комплексных амплитуд гармоник моментов на валах:

\begin{equation}
\begin{aligned}
	\hat{A}_i = A_{0,i} + c_xA_{\alpha,i} + c_yA_{\beta,i},~i=1,\ldots,n
\end{aligned}
\end{equation} 

Для каждого $i$ это трехмерный вектор с компонентами, относящимися к каждому кривошипу.
Результаты идентификации, приведенные в главе 6, показывают, что измерения имеют существенное запаздывание по фазе, зависящие от частоты.
%
%По результатам идентификации для платформы с исходными данными $a = 0,25 м; b = 0,4 м; c = 0,5 м; R = 0,15 м; L = 0,56 м$ были определены следующие параметры модели:
%\begin{equation}
%\begin{aligned}
%	....
%\end{aligned}
%\end{equation}

\section{Динамическая ограничительная характеристика}

Для большинства типовых имитационных эффектов движения, характерных для тренажеров АТ и СТ, компоненты движущих моментов, зависящие от угловых и поступательных скоростей и ускорений, оказывают меньшее влияние на результирующие моменты по сравнению со статическими компонентами. Таким образом, для оценки достижимости положений при планировании траектории в первую очередь необходимо построить ограничительную статическую характеристику, которая будет оказывать максимальное влияние на работоспособность системы. Однако данная статическая характеристика должна уточняться с учетом мощностных характеристик используемых двигателей. В результате чего, ее конечный вид будет определяться суперпозицией двух характеристик (статической и механической двигателей).

\begin{figure}[ht]
    \center{
        %\hfill
        \subcaptionbox[List-of-Figures entry]{\label{img:dyn_restrict_Q800}} {\includegraphics[width=0.45\linewidth]{dyn_restrict_Q800}}%
        \hfill       
        \subcaptionbox{\label{img:dyn_restrict_Q600}} {\includegraphics[width=0.45\linewidth]{dyn_restrict_Q600}}
        \hfill
    }
    \center{
        \subcaptionbox[List-of-Figures entry]{\label{img:dyn_restrict_Q500}} {\includegraphics[width=0.45\linewidth]{dyn_restrict_Q500}}
        \hfill
        \subcaptionbox[List-of-Figures entry]{\label{img:dyn_restrict_Q350}} {\includegraphics[width=0.45\linewidth]{dyn_restrict_Q350}}
        \hfill
    }
    \caption{Динамическая ограничительная характеристика 3DOF СП в декартовых координатах} % Этот текст попадает в названия рисунков в списке рисунков
    \label{img:dyn_restrict_static}
\end{figure}

На рис. 2(а) проиллюстрирована ограничительная характеристика для платформы с параметрами $a = 0,25 \text{м}; b = 0,4 \text{м}; c = 0,5 \text{м}; R = 0,15 \text{м}; L = 0,56 \text{м}$ при установленном ограничении на движущие моменты 100\% от номинального электромагнитного момента используемых мотор-редукторов. 
Рабочая зона представлена в виде проекции параллелепипеда на плоскость «крен-тангаж», допустимые высоты в метрах выделены цветом в соответствии со шкалой в правой части рисунка.  
На рис. 2(б) показана рабочая зона с установленным ограничением 50\% от номинального момента. Как видно по срезам рабочей зоны, некорректный выбор двигателей или превышение допустимой нагрузки на двигатели может существенно ограничить допустимые перемещения платформы даже в статическом режиме. Наличие срезов характеристики объясняется неравномерностью распределения нагрузки на двигатели (все кривошипы работают в одной полуплоскости) и несимметричностью расположения двигателей 1 и 2 по сравнению с двигателем 3 вдоль оси Х.

В ходе экспериментов с динамической моделью была также выявлена существенная зависимость значений моментов от величины подъема центра тяжести относительно подвижной платформы.
